\documentclass[11pt,a4paper]{article}
\usepackage[margin=2.2cm]{geometry}
\usepackage{amsmath,amssymb}
\usepackage{graphicx}
\usepackage{booktabs}
\usepackage{siunitx}
\usepackage{xcolor}
\usepackage{hyperref}
\usepackage{caption}
\usepackage[section]{placeins}
\graphicspath{{figs/}}
\newcommand{\HH}{H/H_0}

\title{Robustness of the SHiP muon shield to parameter tolerances}
\author{}
\date{\today}

\begin{document}
\maketitle

\begin{abstract}
How much do the muon hits on the SHiP sensitive plane change when the 43 free parameters of the
optimised muon shield deviate from their nominal values? A neural surrogate of the per-muon hit
probability (an ensemble of four DeepONets) gives the gradient, the Hessian and candidate
worst-case directions, and every conclusion is checked with the full simulation on $10^9$
muons. The nominal design gives about $830$ hits per $10^9$ muons. If every parameter may be off
by up to $\varepsilon$, the worst design found within these tolerances (searched on the surrogate,
then simulated) multiplies the hits by $1.72\pm0.07$ at $\varepsilon=1\%$, $3.8\pm0.2$ at $2\%$,
$9.0\pm0.3$ at $3\%$, $132\pm5$ at $5\%$ and $3440\pm120$ at $10\%$. A single parameter at
its $\pm10\%$ limit gives at most $\times2.1\pm0.1$ (the field of M7). Random independent
deviations are much milder: harmless at $\pm2\%$, mostly harmless at $\pm5\%$, with a heavy
tail at $\pm10\%$. Bounding the worst case takes sub-percent tolerances on a few parameters,
mostly of M7 and M2: the allocation found for at most $\times2$ (simulated: $\times1.56$) holds
M7's field, M7's length and M2's entrance width to $\pm0.5\%$. Near the nominal design every
harmful direction acts through muons of $30$--$75$\,GeV, while the muons above $150$\,GeV respond
the opposite way.
\end{abstract}

\tableofcontents

%==============================================================================
\section{Setup}
%==============================================================================

\paragraph{The shield.}
The design is the optimised configuration (\texttt{try\_opt}) of seven magnets M1--M7 (M0 is
fixed), with 43 free parameters (Table~\ref{tab:params}): lengths \texttt{Z\_len} and gaps
\texttt{Z\_gap}, the transverse sizes \texttt{dXIn}, \texttt{dXOut}, the inner gaps
\texttt{gapIn}, \texttt{gapOut}, the yoke widths \texttt{dXYokeIn/Out} (all in cm) and the field
$B$ of each magnet (the \texttt{NI} column, the target field in tesla). Fields are uniform within
each magnet block; the cavern is included.

\paragraph{Tolerances.}
Each parameter may vary by $\pm10\%$ of its nominal value:
\begin{equation}
  \phi_i = \phi_{0,i}\,(1 + 0.1\,u_i), \qquad u_i\in[-1,1].
\end{equation}
$u=0$ is the nominal design and $u_i=\pm1$ puts parameter $i$ alone at its limit. A step of
$\|u\|=1$ is ``one tolerance unit'': as large as one parameter at its limit, but spread over all
of them (RMS change $1.5\%$ per parameter). The corners of the box are at $\|u\|=\sqrt{43}=6.6$.
M6 and M7 have negative fields, so a negative change of their $B$ lowers $|B|$.

\paragraph{Muons and figure of merit.}
Muons are drawn from the generative flow of \texttt{generator\_muons} (the physical flux) and
placed on axis at $z=-1$. The figure of merit $H$ is the number of muons crossing the last
sensitive plane per $10^9$ muons. At the nominal design, eight simulator runs on the muons of
the analysis give $803$--$852$ (each $\pm29$; the transport is stochastic) and independent muons
$820\pm14$, a hit rate of $8\times10^{-7}$; the surrogate gives $810$. Every simulator number uses $10^9$ muons, the same ones for all designs, with
Poisson errors (the transport is stochastic, so this holds even on shared muons).

%==============================================================================
\section{The surrogate}
%==============================================================================

\paragraph{Model.}
The surrogate (LCSO) models the hit probability of each muon $x$ for a design $u$,
\begin{equation}
  p(\text{hit}\mid x,u) = \sigma\!\big(b(u)\cdot t(x) + c\big),
\end{equation}
a DeepONet with a branch network $b$ for the design and a trunk network $t$ for the muon
($p_x$, $p_y$, $p_z$, charge). The hits $H(u)$ are the integral of $p$ over the muons, a smooth
function of $u$ with exact derivatives. Four DeepONets (two hidden layers of 128 units, 32
features) are trained with different seeds and their logits averaged; the average is itself a
DeepONet with concatenated features. A seed whose final validation loss is more than $10\%$
above the median of its siblings is replaced by a new one (this happened once: $0.0179$ against
$0.0152$--$0.0156$); the rule uses only the held-out training designs, not the references. Errors on surrogate quantities are the standard error of
the four members: the surrogate's own seed-to-seed uncertainty.

\paragraph{Training data.}
The training data are simulated once, in eight pools (Table~\ref{tab:pools}), and reused for every
epoch; nothing is re-simulated during training. Each pool is a set of designs, each simulated
with its own fresh muons from the flow:
\begin{itemize}
  \item \emph{Designs.} $u=\delta\,s\,r$, with $s$ a scrambled Sobol point in $[-1,1]^{43}$ and
        $r$ uniform in $(0,1)$, for a region of half-width $\delta$ ($\delta=1$: the $\pm10\%$
        box). The radial factor $r$ concentrates the designs towards the nominal one: of the
        13\,947 designs drawn this way, 4475 have every parameter within $\pm3\%$. The
        first 16 designs of each base pool are the nominal design itself ($8\times10^8$ muons
        at $u=0$ in all).
  \item \emph{Far from the nominal design.} The surrogate was weakest there (Table~\ref{tab:val}),
        so two pools were added. The \emph{uniform} pool drops the radial factor: 1024 Sobol
        designs uniform in the box, all with some parameter beyond $\pm8.5\%$ (two crashed the
        simulator and are missing). The \emph{adaptive} pool (212 designs) targets where the
        conclusions rest and where the surrogate was least sure: the worst designs found within
        $\pm1$--$10\%$ and those of the tolerance allocations (Sec.~\ref{sec:tolerance}), each with
        four points on its line from the nominal design and 12 perturbations within its box (24
        for the one the simulator showed mispredicted, at $\pm10\%$), and the 64 of 12\,288 random
        designs on which the members of a first ensemble, trained on the other six pools, disagreed
        most.
  \item \emph{Muons.} Drawn from the flow at a raised \emph{momentum} temperature $T$ (only the
        momentum part of the flow is tempered), which broadens the momentum spectrum and raises
        the hit rate from $8\times10^{-7}$ (the physical flux) to $1.4\times10^{-4}$ ($T=1.5$)
        and $1.8\times10^{-3}$ ($T=2$). This does not bias the model, which learns
        $p(\text{hit}\mid x,u)$ per muon; the temperature only decides where in momentum the
        training data is dense. The $T=1$ pool, at the physical flux, holds only 6900 hits.
  \item \emph{Rows.} Every hit and $0.4\%$ of the misses are kept, a case-control sample; the
        logit is fitted with the offset $\log 250$, so the model remains one of the full
        probability.
  \item \emph{Split.} The designs, not the rows, are split: 13\,663 for training and 1518
        held out for early stopping. The validation references of Sec.~\ref{sec:validation} are
        separate simulations, never used for training.
\end{itemize}
In all, $6.2\times10^{10}$ simulated muons give $3.4\times10^7$ hits and $2.8\times10^8$ rows.
At the physical flux, as many hits would need $4\times10^{13}$ muons, i.e.\ about $85\,000$
designs of $5\times10^8$ muons; the same number of simulated muons would buy only $\sim120$
such designs.

\begin{table}[h]
\centering
\caption{Training pools. Designs are drawn in the region given, pulled towards the nominal
design except in the uniform and adaptive pools (see text); designs that crashed the simulator
are missing.}
\label{tab:pools}
\begin{tabular}{lrrrrr}
\toprule
pool & designs & muons per design & momentum $T$ & region & hits \\
\midrule
main           & 2048    & $10^7$ & 1.5 & $\pm10\%$ & $3.0\times10^6$ \\
small region   & 512     & $10^7$ & 1.5 & $\pm5\%$  & $7.0\times10^5$ \\
large region   & 363     & $10^7$ & 1.5 & $\pm20\%$ & $7.6\times10^5$ \\
cold           & 512     & $10^7$ & 1   & $\pm10\%$ & $6.9\times10^3$ \\
hot            & 512     & $10^7$ & 2   & $\pm10\%$ & $9.2\times10^6$ \\
many designs   & 10\,000 & $10^6$ & 2   & $\pm10\%$ & $1.8\times10^7$ \\
uniform        & 1022    & $10^7$ & 1.5 & $\pm10\%$, uniform & $1.8\times10^6$ \\
adaptive       & 212     & $10^7$ & 1.5 & selected  & $4.0\times10^5$ \\
\midrule
total          & 15\,181 & $6.2\times10^{10}$ & & & $3.4\times10^7$ \\
\bottomrule
\end{tabular}
\end{table}

\paragraph{Integration over the muons.}
With the muons on axis, a muon is its momentum, azimuth and charge, which are independent in the
flow: $(p_z,p_t)$ is the image of a 2-d standard normal latent. The surrogate's hits, gradient,
Hessian and per-momentum splits are therefore computed on a quadrature grid (latent $\times$
azimuth $\times$ charge, App.~\ref{app:quadrature}) instead of by sampling muons: exact to
$10^{-5}$ in seconds, where $10^9$ sampled muons leave $2$--$5\%$ of noise. Only the simulator
uses sampled muons.

\subsection{Validation against the simulator}
\label{sec:validation}

Reference designs were simulated on stored muons, on which the surrogate also predicts: $10^9$
muons of the physical flux (the nominal design four times, 16 designs in the half box
$|u_i|\le\tfrac12$, 16 in the box), and a high-momentum check sample of $4\times10^8$ muons drawn
like the $T=1.5$ pools ($\sim5\times10^4$ hits per design, so precise enough for one-at-a-time
effects: the same designs plus every parameter alone at $\pm10\%$). Table~\ref{tab:val} and
Fig.~\ref{fig:val} compare the surrogate with the simulator on them.

\begin{table}[t]
\centering
\caption{Surrogate vs.\ simulator on the reference designs: median $|\log(\text{pred}/\text{sim})|$
(signed for the nominal design), and the relative vector errors of the 43 one-at-a-time (OAT)
secants $\tfrac12[H(+e_i)-H(-e_i)]$ and curvatures $H(+e_i)+H(-e_i)-2H_0$ on the check sample.
``Noise'': what a perfect model would score, from the simulator's Poisson noise.}
\label{tab:val}
\begin{tabular}{lcc}
\toprule
 & noise & surrogate \\
\midrule
Nominal design                  & $1.7\%$  & $-2.0\%$ \\
Designs in the half box         & $2.1\%$  & $4.1\%$ \\
Designs in the box              & $1.6\%$  & $13\%$ \\
Check sample: nominal design    & $0.2\%$  & $+0.6\%$ \\
Check sample: designs in the box & $0.25\%$ & $3.9\%$ \\
OAT secants (cosine)            & $1.6\%$  & $3.9\%$ (0.999) \\
OAT curvatures (cosine)         & $7.6\%$  & $12\%$ (0.993) \\
\bottomrule
\end{tabular}
\end{table}

\begin{figure}[t]
  \centering
  \includegraphics[width=\linewidth]{study_shield_oat_final.png}\\[1ex]
  \includegraphics[width=0.75\linewidth]{study_shield_calibration_final.png}
  \caption{Top: every parameter alone at $\pm10\%$, secant $\tfrac12[H(+10\%)-H(-10\%)]$ on the
  check sample: simulator (squares) and surrogate (circles). Bottom: predicted vs.\ simulated hits
  of the reference designs.}
  \label{fig:val}
\end{figure}

The one-at-a-time effects are reproduced to $3.9\%$, the designs in the half box to $4\%$. The
weak point is far from the nominal design ($13\%$ in the box):
there the hits come mostly from muons of $30$--$75$\,GeV and below, which are rare in every
training sample, while the high-momentum hits that dominate the training data are predicted to a
few percent. Within one tolerance unit, where the derivatives are read, the walks of
Sec.~\ref{sec:walks} agree with the simulator to $1$--$2\sigma$.

The derivatives are well determined: the four members' gradients at $u=0$ agree to
$\cos=0.974$--$0.992$ and their top Hessian eigenvectors to $0.997$--$0.999$ with the
ensemble's. The second and third eigenvectors are nearly
degenerate and not individually determined. How the architecture and training were chosen is
summarised in App.~\ref{app:study}.

%==============================================================================
\section{Method}
%==============================================================================

From the surrogate we compute
\begin{itemize}
  \item the gradient $g=\nabla_u H$ and Hessian $\nabla^2_u H$ at $u=0$, and the exact
        one-at-a-time change $H(\pm e_i)/H_0-1$ of each parameter at its limit;
  \item the active subspace, the eigenvectors of $C=\mathbb{E}_x[\nabla_u p\,\nabla_u p^\top]$:
        the directions in which the hit probability of individual muons changes most;
  \item candidate worst directions, each of unit length:
    \texttt{ascent} ($v\propto g$),
    \texttt{worst} (maximiser of the quadratic model $g\cdot v+\tfrac12 v^\top\nabla^2H\,v$ on
    $\|v\|=1$),
    \texttt{hess1} (top Hessian eigenvector), \texttt{act1}, \texttt{act2} (top active
    directions), \texttt{corner} ($v\propto\operatorname{sign} g$: every parameter moved by the
    same amount) and \texttt{random};
  \item the worst corner of the box, by projected-gradient ascent (\texttt{pga});
  \item 4096 random designs in the box;
  \item the worst design for a given tolerance and a tolerance allocation
        (Sec.~\ref{sec:tolerance}).
\end{itemize}
Each direction is then walked, $u=t\,v$ for $t\in\{-1,-\tfrac12,\tfrac12,1\}$ (for \texttt{pga},
up to the box face), and simulated.

%==============================================================================
\section{Results}
%==============================================================================

\subsection{Single parameters}

32 of the 43 one-at-a-time effects are resolved beyond three times the members' scatter
(Fig.~\ref{fig:sensitivity}, Table~\ref{tab:oat}). The most sensitive parameter is the
\emph{field of M7}: lowering $|B|$ by $10\%$ doubles the hits. Next come the entrance width of M2
(narrowed by $10\%$: also $\times2$), the length, yoke widths and exit width of M7, the inner yoke
of M3, the field of M2 and the length of M6. Most responses are one-sided: one side of the
tolerance hardly matters, the other is harmful. A few (the lengths of M2 and M3, the field of M3)
are harmful on both sides, i.e.\ the nominal value sits at their minimum. Per magnet
(Fig.~\ref{fig:magnets}), M7 carries by far the largest share of the gradient ($77\%$ of its
norm), then M2 ($42\%$), M3 ($32\%$) and M6 ($27\%$).

\paragraph{Simulated.} The eight largest were simulated at their worse limit
(Table~\ref{tab:concrete}): all eight rise significantly, six within $1$--$2\sigma$ of the
surrogate; the exit width of M7 rises more than predicted ($\times1.85\pm0.08$ vs.\ $1.41$) and
the inner yoke of M3 less ($1.26\pm0.06$ vs.\ $1.45$). Their scans were also simulated at
$\pm5\%$ and $\pm10\%$ (Fig.~\ref{fig:scans}): 25 of the 32 points agree within $2\sigma$ (14 of
16 at $\pm5\%$). The misses are the downstream geometry of M7 narrowed, where the surrogate is too
smooth (exit width at $-5\%$ and $-10\%$: $1.36$ and $1.91$ simulated vs.\ $1.15$ and $1.41$; inner
yoke at $-10\%$: $1.25$ vs.\ $1.02$), and the entrance width of M2 widened by $10\%$, harmless in
the simulator ($1.01\pm0.05$) where the surrogate predicts a small rise ($1.17$). The inner yoke of
M7 is harmful on both sides in the simulator.

\begin{table}[t]
\centering
\caption{The most sensitive parameters (surrogate): first-order sensitivity $g_i/H_0$ per full
tolerance, and the exact change at $\pm10\%$ (errors: the members' standard error).}
\label{tab:oat}
\begin{tabular}{lrrrr}
\toprule
parameter & nominal & $g_i/H_0$ & $H/H_0-1$ at $+10\%$ & at $-10\%$ \\
\midrule
M7:$B$                 & $-1.9$ T & $-0.494$ & $-0.26\pm0.01$ & $+1.03\pm0.02$ \\
M2:\texttt{dXIn}       & 68.04 cm & $-0.384$ & $+0.17\pm0.02$ & $+0.99\pm0.03$ \\
M7:\texttt{Z\_len}     & 193.1 cm & $-0.389$ & $-0.23\pm0.01$ & $+0.69\pm0.03$ \\
M7:\texttt{dXYokeIn}   & 124.8 cm & $+0.403$ & $+0.47\pm0.03$ & $+0.02\pm0.03$ \\
M3:\texttt{dXYokeIn}   & 85.88 cm & $-0.315$ & $-0.11\pm0.01$ & $+0.45\pm0.04$ \\
M7:\texttt{dXOut}      & 67.81 cm & $-0.221$ & $-0.11\pm0.02$ & $+0.41\pm0.03$ \\
M2:$B$                 & 1.9 T    & $-0.209$ & $-0.14\pm0.02$ & $+0.34\pm0.02$ \\
M6:\texttt{Z\_len}     & 168.0 cm & $-0.240$ & $-0.19\pm0.01$ & $+0.30\pm0.01$ \\
M7:\texttt{dXYokeOut}  & 94.92 cm & $+0.251$ & $+0.27\pm0.02$ & $-0.09\pm0.01$ \\
M2:\texttt{Z\_len}     & 224.9 cm & $-0.040$ & $+0.10\pm0.03$ & $+0.25\pm0.02$ \\
M4:\texttt{Z\_len}     & 224.0 cm & $-0.127$ & $-0.07\pm0.02$ & $+0.25\pm0.01$ \\
M3:$B$                 & 1.9 T    & $+0.059$ & $+0.24\pm0.03$ & $+0.13\pm0.02$ \\
M3:\texttt{Z\_len}     & 225.0 cm & $-0.013$ & $+0.14\pm0.02$ & $+0.22\pm0.02$ \\
M1:$B$                 & 1.9 T    & $-0.149$ & $-0.10\pm0.01$ & $+0.22\pm0.01$ \\
M3:\texttt{dXIn}       & 64.81 cm & $+0.110$ & $+0.21\pm0.03$ & $-0.02\pm0.01$ \\
\bottomrule
\end{tabular}
\end{table}

\begin{figure}[p]
  \centering
  \includegraphics[width=\linewidth]{robustness_sensitivity.png}
  \caption{Sensitivity of each parameter. Left: gradient per full tolerance. Right: exact
  change with the parameter at $+10\%$ (orange) and $-10\%$ (blue).}
  \label{fig:sensitivity}
\end{figure}

\begin{figure}[t]
  \centering
  \includegraphics[width=0.8\linewidth]{robustness_magnets.png}
  \caption{Per magnet: share of the gradient norm and largest single-parameter effect.}
  \label{fig:magnets}
\end{figure}

\begin{figure}[t]
  \centering
  \includegraphics[width=\linewidth]{robustness_scans.png}
  \caption{One-at-a-time scans of the eight most sensitive parameters: surrogate (blue) and
  simulator at $\pm5\%$ and $\pm10\%$ (orange).}
  \label{fig:scans}
\end{figure}

\subsection{Curvature and active subspace}

One Hessian eigenvalue dominates (Fig.~\ref{fig:spectra}): $3.41\pm0.13\,H_0$ per unit$^2$, then
$1.01$, $0.86$ and $\lesssim0.17\,H_0$; 14 are negative, none below $-0.18\,H_0$. Along the
gradient the second-order term already exceeds the first-order one at one tolerance unit
($1.19\,H_0$ vs.\ $1.07\,H_0$): the response is strongly non-linear at that scale.

The active subspace is effectively one-dimensional ($98.8\%$ of $\mathbb{E}[\|\nabla_u p\|^2]$).
Its direction \texttt{act1} shortens M1--M4 and lowers their fields together, i.e.\ it changes the
upstream bending power (Table~\ref{tab:dirs}). It is almost orthogonal to the gradient
($\cos=0.06$): individual muons are deflected most along it, but the changes largely cancel in
the total.

\begin{figure}[t]
  \centering
  \includegraphics[width=0.9\linewidth]{robustness_spectra.png}
  \caption{Left: Hessian eigenvalues in units of $H_0$. Right: active-subspace eigenvalues
  (blue), their cumulative share (red) and the fraction of the gradient they contain (grey).}
  \label{fig:spectra}
\end{figure}

\subsection{Worst directions}
\label{sec:walks}

\begin{table}[t]
\centering
\caption{Hits one tolerance unit along each direction, surrogate and simulator (nominal: $843$
simulated hits). ``max / RMS'': the largest and the RMS change over the 43 parameters.}
\label{tab:walks}
\begin{tabular}{lcccrr}
\toprule
direction & max / RMS change & surrogate $\HH$ & simulator $\HH$ & sim.\ hits & at $t=-1$ (sim.) \\
\midrule
\texttt{worst}  & 4.3\% / 1.5\% & $5.20\pm0.17$ & $4.95\pm0.19$ & 4174 & $1.00\pm0.05$ \\
\texttt{hess1}  & 4.2\% / 1.5\% & $5.26\pm0.18$ & $5.19\pm0.20$ & 4377 & $1.31\pm0.06$ \\
\texttt{ascent} & 4.6\% / 1.5\% & $3.86\pm0.06$ & $4.02\pm0.15$ & 3389 & $0.60\pm0.03$ \\
\texttt{act1}   & 3.9\% / 1.5\% & $2.56\pm0.07$ & $2.44\pm0.10$ & 2059 & $1.50\pm0.07$ \\
\texttt{corner} & 1.5\% / 1.5\% & $2.15\pm0.04$ & $2.01\pm0.08$ & 1690 & $0.67\pm0.04$ \\
\texttt{act2}   & 9.4\% / 1.5\% & $1.20\pm0.03$ & $1.39\pm0.06$ & 1174 & $0.87\pm0.04$ \\
\texttt{random} & ---  / 1.5\% & $0.96\pm0.00$ & $0.87\pm0.04$ &  731 & $1.10\pm0.05$ \\
\bottomrule
\end{tabular}
\end{table}

\paragraph{Small coordinated changes are far more harmful than a large single one.}
Moving every parameter by at most $4.3\%$ (RMS $1.5\%$) along \texttt{worst} multiplies the hits
by $4.95\pm0.19$ (Table~\ref{tab:walks}). The simplest such pattern, \texttt{corner}, moves
\emph{every} parameter by exactly $1.5\%$ with the sign of its gradient and doubles the hits. The
worst single parameter needs a $10\%$ error for a similar effect. A random direction of the same
size changes little ($0.87\pm0.04$): the damage comes from the pattern of signs, not the size of
the step.

\paragraph{The surrogate gets both the directions and the sizes.}
Its two best directions (\texttt{worst} and \texttt{hess1}, $\cos=0.98$) are the two the
simulator ranks highest, then \texttt{ascent}. Within one tolerance unit it agrees with the
simulator to $1$--$2\sigma$ almost everywhere (Fig.~\ref{fig:calibration}); the exceptions are
\texttt{act2} at $t=1$ ($1.20$ vs.\ $1.39\pm0.06$) and the $t=-1$ ends of \texttt{ascent} and
\texttt{worst}, where it is $50\%$ and $17\%$ high.

\paragraph{The quadratic model is not enough.}
Built from the gradient and Hessian, it underestimates every harmful direction ($3.6$ instead of
$4.95$ along \texttt{worst}, Fig.~\ref{fig:walks}): the hits grow roughly exponentially. The
derivatives give the directions; magnitudes come from the surrogate and the simulator.

\paragraph{The active subspace is a direction of curvature, not of slope.}
\texttt{act1} gives $\times2.44$ at $t=+1$ and also $\times1.50$ at $t=-1$: shortening or
lengthening magnets 1--4 together, with their fields, is harmful either way, and the nominal
design sits at the balance point (Sec.~\ref{sec:momentum}). In the plane of \texttt{worst} and
\texttt{act1} (Fig.~\ref{fig:plane}) the nominal design lies in a shallow valley.

\begin{table}[t]
\centering
\caption{What each direction changes at one tolerance unit: its largest parameter changes,
relative and in physical units (cm, or T for $B$).}
\label{tab:dirs}
\small
\begin{tabular}{l p{0.55\linewidth} p{0.25\linewidth}}
\toprule
direction & largest changes & in words \\
\midrule
\texttt{worst}  & M2:\texttt{dXIn} $-4.3\%$ ($-2.9$), M7:$B$ $-4.2\%$ ($|B|$ $-0.08$ T),
                  M7:\texttt{Z\_len} $-3.5\%$ ($-6.8$), M7:\texttt{dXYokeIn} $+2.9\%$ ($+3.7$),
                  M3:\texttt{dXYokeIn} $-2.9\%$ ($-2.5$), M3:$B$ $+2.3\%$ ($+0.04$ T),
                  M7:\texttt{dXOut} $-2.3\%$ ($-1.5$)
                & weaker and shorter M7 with wider yokes, narrower M2 entrance \\
\midrule
\texttt{hess1}  & M2:\texttt{dXIn} $-4.2\%$ ($-2.9$), M7:$B$ $-3.7\%$ ($|B|$ $-0.07$ T),
                  M7:\texttt{Z\_len} $-3.2\%$ ($-6.3$), M3:$B$ $+2.9\%$ ($+0.06$ T),
                  M3:\texttt{dXYokeIn} $-2.6\%$ ($-2.2$), M2:\texttt{Z\_len} $+2.5\%$ ($+5.6$),
                  M3:\texttt{Z\_len} $+2.4\%$ ($+5.3$)
                & as \texttt{worst} \\
\midrule
\texttt{ascent} & M7:$B$ $-4.6\%$ ($|B|$ $-0.09$ T), M7:\texttt{dXYokeIn} $+3.8\%$ ($+4.7$),
                  M7:\texttt{Z\_len} $-3.6\%$ ($-7.0$), M2:\texttt{dXIn} $-3.6\%$ ($-2.4$),
                  M3:\texttt{dXYokeIn} $-2.9\%$ ($-2.5$), M7:\texttt{dXYokeOut} $+2.3\%$ ($+2.2$),
                  M6:\texttt{Z\_len} $-2.2\%$ ($-3.8$)
                & the gradient: M7 weaker and shorter \\
\midrule
\texttt{corner} & all 43 parameters by exactly $\pm1.5\%$, each with the sign of its gradient
                & the simplest uniform pattern \\
\midrule
\texttt{act1}   & M2, M3, M4:\texttt{Z\_len} $-3.9\%$ each ($-8.7$ to $-8.8$), M2:$B$ $-3.5\%$,
                  M3:$B$ $-3.1\%$, M4:$B$ $-2.8\%$ ($-0.05$ to $-0.07$ T), M1:\texttt{Z\_len}
                  $-2.7\%$ ($-4.0$)
                & magnets 1--4 shorter and weaker: less bending upstream \\
\midrule
\texttt{act2}   & M7:\texttt{dXOut} $-9.4\%$ ($-6.4$), M7:\texttt{Z\_len} $+2.5\%$ ($+4.8$), the
                  rest at most $1\%$
                & the exit width of M7 \\
\bottomrule
\end{tabular}
\end{table}

\begin{figure}[p]
  \centering
  \includegraphics[width=\linewidth]{robustness_walks.png}
  \caption{Walks along each direction: surrogate (blue; error bars: the members), simulator
  (orange) and the quadratic model (dotted).}
  \label{fig:walks}
\end{figure}

\begin{figure}[t]
  \centering
  \includegraphics[width=0.48\linewidth]{robustness_summary.png}\hfill
  \includegraphics[width=0.42\linewidth]{robustness_calibration.png}
  \caption{Left: $\HH$ at one tolerance unit along each direction. Right: surrogate vs.\
  simulator for all 33 walked designs; the outliers are the box designs of Sec.~\ref{sec:box}.}
  \label{fig:calibration}
\end{figure}

\begin{figure}[t]
  \centering
  \includegraphics[width=0.6\linewidth]{robustness_plane.png}
  \caption{The surrogate in the plane of \texttt{worst} ($x$) and the part of \texttt{act1}
  orthogonal to it ($y$); markers: simulated designs with their $\HH$.}
  \label{fig:plane}
\end{figure}

\subsection{Concrete designs along the worst direction}
\label{sec:concrete}

Table~\ref{tab:concrete} lists concrete designs, all simulated. ``Worst, max $\varepsilon$'' is
the worst direction scaled so that no parameter moves by more than $\varepsilon$. It lies within
the $\pm\varepsilon$ tolerances, so it bounds their worst case from below; most of its
parameters stay well inside them, and the worst design within the tolerances is much worse
(Sec.~\ref{sec:tolerance}).

\begin{table}[t]
\centering
\caption{Concrete designs, simulated (nominal $847\pm29$ hits) and predicted. At
$\varepsilon=10\%$ the worst direction changes M2:\texttt{dXIn} by $-6.8$ cm, $|B|$ of M7 from
$1.9$ to $1.71$ T, M7:\texttt{Z\_len} by $-15.9$ cm, M7:\texttt{dXYokeIn} by $+8.5$ cm,
M3:\texttt{dXYokeIn} by $-5.7$ cm, M3:$B$ by $+0.10$ T, M7:\texttt{dXOut} by $-3.6$ cm,
M7:\texttt{dXYokeOut} by $+4.1$ cm, M3:\texttt{dXIn} by $+2.6$ cm, M2:\texttt{Z\_len} by $+8.9$ cm,
and 33 more parameters by less than $4\%$.}
\label{tab:concrete}
\small
\begin{tabular}{llrcc}
\toprule
design & what changes & hits & sim.\ $\HH$ & surr.\ $\HH$ \\
\midrule
worst, max 1\%   & all 43, RMS 0.35\% & 1047  & $1.24\pm0.06$ & 1.32 \\
worst, max 2\%   & all 43, RMS 0.7\% & 1499  & $1.77\pm0.08$ & 1.93 \\
worst, max 3\%   & all 43, RMS 1.05\% & 2276  & $2.69\pm0.11$ & 2.95 \\
worst, max 5\%   & all 43, RMS 1.75\% & 5741  & $6.78\pm0.25$ & 7.05 \\
worst, max 10\%  & all 43, RMS 3.5\% & 61579 & $72.7\pm2.5$  & 54.9 \\
opposite, max 5\% & all 43, signs reversed & 999 & $1.18\pm0.06$ & 1.40 \\
opposite, max 10\% & all 43, signs reversed & 4409 & $5.21\pm0.20$ & 6.46 \\
\midrule
worst, top 1 only  & M2:\texttt{dXIn} $-10\%$ ($-6.8$ cm) & 1594 & $1.88\pm0.08$ & 1.99 \\
worst, top 3 only  & + M7:$B$ $-10\%$, M7:\texttt{Z\_len} $-8\%$ & 5824 & $6.88\pm0.25$ & 6.47 \\
worst, top 5 only  & + two yoke widths ($\pm6.8\%$) & 10937 & $12.9\pm0.5$ & 11.1 \\
worst, top 10 only & the ten largest changes & 29146 & $34.4\pm1.2$ & 28.8 \\
\midrule
single & M7:$B$ $-10\%$ ($|B|$ $1.9\to1.71$ T) & 1790 & $2.11\pm0.09$ & 2.03 \\
single & M2:\texttt{dXIn} $-10\%$ ($-6.8$ cm)   & 1656 & $1.96\pm0.08$ & 1.99 \\
single & M7:\texttt{dXOut} $-10\%$ ($-6.8$ cm)  & 1563 & $1.85\pm0.08$ & 1.41 \\
single & M7:\texttt{Z\_len} $-10\%$ ($-19.3$ cm) & 1395 & $1.65\pm0.07$ & 1.69 \\
single & M7:\texttt{dXYokeIn} $+10\%$ ($+12.5$ cm) & 1315 & $1.55\pm0.07$ & 1.47 \\
single & M2:$B$ $-10\%$ ($1.9\to1.71$ T) & 1071 & $1.26\pm0.06$ & 1.34 \\
single & M3:\texttt{dXYokeIn} $-10\%$ ($-8.6$ cm) & 1070 & $1.26\pm0.06$ & 1.45 \\
single & M6:\texttt{Z\_len} $-10\%$ ($-16.8$ cm) & 1040 & $1.23\pm0.06$ & 1.30 \\
\bottomrule
\end{tabular}
\end{table}

\paragraph{Exponential growth along the direction.}
Along the worst direction (Fig.~\ref{fig:scaled}) the hits grow by $\times1.24$ at a maximum
change of $1\%$, $\times1.8$ at $2\%$, $\times2.7$ at $3\%$, $\times6.8$ at $5\%$ and $\times73$ at
$10\%$: a factor $e$ per $\sim2.3\%$ of tolerance. The surrogate follows to within $10\%$ up to
$5\%$ and is $25\%$ low at $10\%$. The opposite direction is harmless at $5\%$ ($\times1.2$) but
not at $10\%$ ($\times5.2$), through other muons (Sec.~\ref{sec:momentum}).

\paragraph{The rise is collective.}
The largest component alone (M2:\texttt{dXIn} $-10\%$) gives $\times1.9$; the top three (adding a
weaker and shorter M7) $\times6.9$; the top five $\times13$; the top ten $\times34$; all 43
$\times73$ (Fig.~\ref{fig:attribution}). Each parameter contributes a moderate factor, and the
factors multiply.

\begin{figure}[t]
  \centering
  \includegraphics[width=0.7\linewidth]{worst_direction_scaled.png}
  \caption{Hits along the worst direction against the largest change of any parameter:
  simulator and surrogate.}
  \label{fig:scaled}
\end{figure}

\begin{figure}[t]
  \centering
  \includegraphics[width=0.8\linewidth]{worst_parameters.png}
  \caption{The worst direction at a maximum change of $10\%$, truncated to its largest 1, 3, 5,
  10 changes, and single parameters at their worse limit.}
  \label{fig:attribution}
\end{figure}

\FloatBarrier
\subsection{The worst case for a given tolerance}
\label{sec:tolerance}
\label{sec:box}

\paragraph{Method.} If every parameter may be off by up to $\varepsilon$, the designs allowed
form the box $|u_i|\le\varepsilon/10\%$. Its worst design is found on the surrogate, by
projected-gradient ascent of $\log H$ over the box from 16 starting points (the sign of the
gradient, the scaled worst direction and random corners; 300 Adam steps each, all batched,
seconds per box), and then simulated on the $10^9$ muons of the walks. The simulated values are
real designs within the tolerances, so the true worst case is at least as bad; the search is as
good as the surrogate, which is accurate there up to $\pm3\%$ and overestimates beyond. The worst
direction of Sec.~\ref{sec:walks} maximises the hits on a sphere ($\|u\|=1$); on a box the
maximum is at a corner, every parameter at a limit with the harmful sign, so scaling the
direction down underestimates it.

\begin{table}[!htb]
\centering
\caption{The worst design within $\pm\varepsilon$ on every parameter: found on the surrogate,
simulated ($10^9$ muons; nominal $849\pm21$, two runs). Last column: the worst direction scaled
to the same $\varepsilon$ (Table~\ref{tab:concrete}).}
\label{tab:tolerance}
\begin{tabular}{lrrccc}
\toprule
$\varepsilon$ & at a limit & sim.\ hits & sim.\ $\HH$ & surr.\ $\HH$ & scaled direction (sim.) \\
\midrule
$1\%$  & 43 & 1462      & $1.72\pm0.07$ & $1.79\pm0.04$ & $1.24\pm0.06$ \\
$2\%$  & 43 & 3207      & $3.78\pm0.15$ & $3.91\pm0.15$ & $1.77\pm0.08$ \\
$3\%$  & 43 & 7644      & $9.0\pm0.3$   & $9.5\pm0.4$   & $2.69\pm0.11$ \\
$5\%$  & 39 & 112\,445  & $132\pm5$     & $236\pm18$    & $6.78\pm0.25$ \\
$10\%$ & 39 & 2\,917\,029 & $3440\pm120$ & $(1.1\pm0.7)\times10^5$ & $72.7\pm2.5$ \\
\bottomrule
\end{tabular}
\end{table}

\begin{figure}[!htb]
  \centering
  \includegraphics[width=0.75\linewidth]{tolerance_boxes.png}
  \caption{Worst case for a given tolerance: the worst design within $\pm\varepsilon$ (filled;
  surrogate and simulator) and the worst direction scaled to $\varepsilon$ (open).}
  \label{fig:tolerance}
\end{figure}

\paragraph{Results} (Table~\ref{tab:tolerance}, Fig.~\ref{fig:tolerance}). Already at $\pm1\%$
the worst design gives $\times1.7$, at $\pm2\%$ $\times3.8$ and at $\pm3\%$ $\times9.0$, $1.4$ to
three times the scaled direction. Up to $\pm3\%$ these designs act through the known mechanism:
the largest contributions are those of the worst direction (M7 weaker and shorter, M2's entrance
narrower, M7's exit and M3's inner yoke narrower), but every parameter is at its limit
($\cos\approx0.62$ with the worst direction), and the surrogate predicts them to $3$--$5\%$. At
$\pm5\%$ and $\pm10\%$ the search finds another mechanism, M7's yokes narrower with M2's entrance
and M3's inner yoke wider ($\cos=-0.26$ and $-0.12$ with the worst direction), which the local
derivatives do not show; the simulated scans point the same way (M7's inner yoke harmful on both
sides, Fig.~\ref{fig:scans}). There the surrogate overestimates, by $\times1.8$ at $\pm5\%$ and
$\times30$ at $\pm10\%$, but the designs it finds are real: $\times132$ at $\pm5\%$ and $\times3440$
at $\pm10\%$. The cruder box search of the walks (\texttt{pga}, 4 starts, 200 steps) ends at a
corner with $\times1070\pm40$, almost all from muons of $40$--$60$\,GeV; halfway to it (every
parameter at $5\%$) the simulator gives $\times34\pm1$, and its opposite corner $\times28\pm1$.
Which muons make the hits of the $\pm5\%$ and $\pm10\%$ worst designs was not measured.

\subsection{Tolerance allocation}
\label{sec:allocation}

\paragraph{Method.} Which tolerances keep the worst case below a target? Greedily: start with
every parameter at $\pm10\%$; find the box's worst design $u^*$ as above; tighten the two
parameters that contribute most to it (the largest first-order shares of $\log H$,
$u^*_i\,\partial\log H/\partial u_i$) by a factor $0.8$, down to a floor of $\pm0.5\%$; repeat
until the worst case is below the target. The first box that meets each target is its
allocation, and its worst design is simulated. This is a heuristic, not an optimum: other
allocations may loosen some parameters for the same target.

\begin{table}[!htb]
\centering
\caption{Tolerance allocations (greedy, surrogate) for a worst case of at most $3$, $2$ and $1.5$
times the nominal hits: the sixteen tightest parameters ($\pm\%$), and a summary of all 43. The
worst design of each allocation is simulated in the last row.}
\label{tab:allocation}
\small
\begin{tabular}{lrccc}
\toprule
parameter & nominal & $\le3H_0$ & $\le2H_0$ & $\le1.5H_0$ \\
\midrule
M2:\texttt{dXIn}         & 68.04 cm & 0.55 & 0.50 & 0.50 \\
M7:$B$                   & 1.9 T    & 0.50 & 0.50 & 0.50 \\
M2:\texttt{dXIn}         & 68.04 cm & 0.55 & 0.50 & 0.50 \\
M7:\texttt{Z\_len}       & 193.1 cm & 0.55 & 0.50 & 0.50 \\
M7:\texttt{dXYokeIn}     & 124.8 cm & 1.07 & 0.50 & 0.50 \\
M7:\texttt{dXOut}        & 67.81 cm & 0.69 & 0.55 & 0.50 \\
M3:\texttt{dXYokeIn}     & 85.88 cm & 1.07 & 0.55 & 0.50 \\
M3:$B$                   & 1.9 T    & 1.07 & 0.69 & 0.50 \\
M6:\texttt{Z\_len}       & 168.0 cm & 1.34 & 0.69 & 0.50 \\
M3:\texttt{dXIn}         & 64.81 cm & 1.34 & 0.86 & 0.50 \\
M7:\texttt{dXYokeOut}    & 94.92 cm & 1.34 & 0.86 & 0.50 \\
M4:\texttt{dXIn}         & 39.02 cm & 2.10 & 1.07 & 0.55 \\
M4:$B$                   & 1.9 T    & 2.10 & 1.07 & 0.55 \\
M2:\texttt{dXYokeIn}     & 70.62 cm & 2.62 & 1.34 & 0.55 \\
M7:\texttt{dXIn}         & 37.95 cm & 1.68 & 1.07 & 0.69 \\
M2:$B$                   & 1.9 T    & 2.10 & 1.34 & 0.69 \\
M3:\texttt{Z\_len}       & 225.0 cm & 1.34 & 0.86 & 0.86 \\
\midrule
parameters below $\pm1\%$ / $\pm2\%$ / $\pm5\%$ & & 4 / 14 / 30 & 11 / 18 / 36 & 20 / 27 / 41 \\
median tolerance & & $5.1\%$ & $2.6\%$ & $1.1\%$ \\
\midrule
worst design: surrogate & & $2.97\pm0.26$ & $1.99\pm0.06$ & $1.50\pm0.04$ \\
worst design: simulator & & $2.44\pm0.10$ & $1.56\pm0.07$ & $1.40\pm0.06$ \\
\bottomrule
\end{tabular}
\end{table}

\begin{figure}[!htbp]
  \centering
  \includegraphics[width=\linewidth]{tolerance_allocation.png}
  \caption{Left: the tolerance allocations for a worst case of at most $3$, $2$ and $1.5\,H_0$
  (surrogate; simulated worst designs in the legend). Right: the greedy path, the worst case of
  the box against the number of parameters tightened.}
  \label{fig:allocation}
\end{figure}

\paragraph{Results} (Table~\ref{tab:allocation}, Fig.~\ref{fig:allocation}). A guarantee against
the worst case is demanding. For at most $\times3$, the allocation holds M7's field to $\pm0.5\%$
($0.01$ T), M2's entrance width and M7's length to $\pm0.55\%$ ($0.37$ and $1.1$ cm) and M7's exit
width to $\pm0.7\%$; 14 parameters are below $\pm2\%$, two may stay at $\pm10\%$, and the median
tolerance is $5\%$. For at most $\times2$ it puts eleven parameters below $\pm1\%$, and for
$\times1.5$ twenty, ten of them at the $0.5\%$ floor; with every parameter at $\pm0.5\%$ the worst
case is still $\times1.29$. The order follows the sensitivities: M7 (field, length, exit width,
yokes), M2's entrance and M3 (inner yoke, field, entrance) first, the gaps between magnets last.
The simulated worst designs stay below every target ($\times2.44$, $1.56$, $1.40$): at these small
tolerances the surrogate is conservative.

\subsection{Random tolerances}

Designs with every parameter independently uniform in $\pm2\%$, $\pm5\%$ and $\pm10\%$: eight of
each simulated, and 4096 per level from the surrogate (Fig.~\ref{fig:random}).
\begin{center}
\begin{tabular}{lcccc}
\toprule
tolerance & simulated $\HH$ (8 designs) & surr.\ median & surr.\ 5--95\% & $P(H>H_0)$ \\
\midrule
$\pm2\%$  & 0.84, 0.87, 0.89, 0.92, 0.93, 0.94, 0.99, 1.13 & 1.03 & 0.85--1.27 & $58\%$ \\
$\pm5\%$  & 0.93, 0.95, 1.06, 1.07, 1.23, 1.39, 1.52, 4.04 & 1.20 & 0.78--2.16 & $74\%$ \\
$\pm10\%$ & 0.69, 0.83, 1.17, 1.23, 1.84, 3.30, 5.00, 7.43 & 2.08 & 0.82--8.98 & $90\%$ \\
\bottomrule
\end{tabular}
\end{center}
At $\pm2\%$ the hits stay within about $\pm20\%$. At $\pm5\%$ most designs stay within $\pm50\%$,
but one of the eight simulated quadruples the hits. At $\pm10\%$ the median design is twice as
bad and the tail reaches $\times7.4$ among the simulated designs ($\times29$ at the surrogate's
99th percentile). The surrogate agrees with the simulator design by design (median ratio
$0.88$--$1.02$).

\begin{figure}[t]
  \centering
  \includegraphics[width=\linewidth]{random_tolerances.png}
  \caption{Random independent tolerances: surrogate distribution of 4096 designs per level
  (histogram), and the eight simulated designs (solid: simulator, dotted: surrogate).}
  \label{fig:random}
\end{figure}

\subsection{Which muons respond}
\label{sec:momentum}

The muons are split into bins of $p_z$ (and separately $p_t$ and charge), and each bin's hits
$H_b$ get their own gradient, one-at-a-time effects and worst direction (Fig.~\ref{fig:pz_overview}).
Below $75$\,GeV the per-bin magnitudes carry the surrogate's larger uncertainty there
(Sec.~\ref{sec:validation}), so the directions and signs are the robust part; the simulated hits
per $p_z$ range (Fig.~\ref{fig:simpz}) confirm them.

\paragraph{Three populations.}
Half of the muons have $p_z<30$\,GeV, but at the nominal design the hits come from two groups
(Table~\ref{tab:pz}; simulated shares in parentheses): \emph{intermediate} momenta,
$30$--$75$\,GeV, with $66\%$ of the hits ($66\%$; $43\%$ at $40$--$50$\,GeV alone) at a hit rate
of $10^{-6}$; and \emph{high} momenta above $125$\,GeV, $1.2\%$ of the muons but $26\%$ of the hits
($24\%$ above $150$\,GeV), with hit rates up to $10^{-3}$ above $300$\,GeV: the muons the shield
does not deflect enough. A third group, \emph{soft} muons of $20$--$30$\,GeV, gives only $4\%$ of
the hits but responds strongly to some changes.

\begin{table}[t]
\centering
\caption{Per $p_z$ bin at the nominal design (surrogate): share of the muons and of the hits, hit
rate, relative sensitivity $|\nabla H_b|/H_b$ per tolerance unit, and the cosine between the
bin's gradient and that of the total.}
\label{tab:pz}
\begin{tabular}{lrrrrr}
\toprule
$p_z$ [GeV] & muons & share of hits & hit rate & $|\nabla H_b|/H_b$ & $\cos(g_b, g)$ \\
\midrule
10--20   & $1.7\%$   & $0.0\%$  & $1.2\times10^{-8}$ & 1.79 & $-0.17$ \\
20--30   & $48.4\%$  & $3.9\%$  & $6.6\times10^{-8}$ & 2.44 & $-0.42$ \\
30--45   & $28.8\%$  & $42.6\%$ & $1.2\times10^{-6}$ & 1.65 & $+0.94$ \\
45--75   & $15.3\%$  & $23.3\%$ & $1.2\times10^{-6}$ & 2.13 & $+0.83$ \\
75--125  & $4.7\%$   & $4.2\%$  & $7.2\times10^{-7}$ & 0.38 & $+0.40$ \\
125--175 & $0.89\%$  & $2.7\%$  & $2.5\times10^{-6}$ & 0.43 & $+0.08$ \\
175--225 & $0.23\%$  & $3.0\%$  & $1.1\times10^{-5}$ & 1.28 & $+0.07$ \\
225--300 & $0.08\%$  & $10.5\%$ & $1.1\times10^{-4}$ & 1.23 & $+0.07$ \\
$>300$   & $0.006\%$ & $9.7\%$  & $1.3\times10^{-3}$ & 2.96 & $+0.06$ \\
\bottomrule
\end{tabular}
\end{table}

\paragraph{The populations want opposite changes.}
The worst directions of the $30$--$45$ and $45$--$75$\,GeV bins agree ($\cos=0.95$), as do all
bins above $125$\,GeV ($0.89$--$1.00$), but the two groups are anti-aligned ($\cos=-0.24$ to
$-0.47$), and the soft bins are anti-aligned with the intermediate ones too (down to $-0.61$,
Fig.~\ref{fig:pz_alignment}). The total gradient, and every harmful direction of
Table~\ref{tab:walks}, follows the intermediate group ($\cos=0.83$--$0.94$). Simulated
(Fig.~\ref{fig:simpz}):
\begin{itemize}
  \item worst direction at $10\%$: the $40$--$50$\,GeV hits go from $363$ to $44\,637$ and the
        $50$--$60$\,GeV ones from $67$ to $14\,346$, while those above $150$\,GeV \emph{fall}
        from $202$ to $45$ and those below $30$\,GeV from $37$ to $6$;
  \item opposite direction: the $40$--$50$\,GeV hits fall from $363$ to $71$, but those of
        $20$--$30$\,GeV rise from $37$ to $2353$ and those above $150$\,GeV from $202$ to $1286$;
  \item the box corner of the walks: only $30$--$75$\,GeV ($562\to903\,046$), the high-momentum
        hits fall ($202\to42$);
  \item random $\pm10\%$ designs can open any group: one raises the $40$--$50$\,GeV hits from $363$
        to $3060$, another those above $150$\,GeV from $202$ to $1033$, and the worst of the eight
        ($\times7.4$) both the soft ($37\to1967$) and the high-momentum ($202\to3296$) ones.
\end{itemize}

\paragraph{Which parameters act on which muons} (Fig.~\ref{fig:pz_sensitivity}).
The high-momentum hits are controlled by the \emph{upstream} magnets: the lengths and fields of
M1--M4, each changing the hits above $175$\,GeV by $40$--$200\%$ at its worse limit. The
intermediate hits are controlled by the entrance width of M2, the yokes of M3 and M7, and above
all the \emph{downstream} magnet M7 (its field: $+150$--$170\%$ of the $30$--$75$\,GeV hits at the
worse limit). The soft hits respond most to the entrance widths of M2 and M3 and the yoke of M3
(up to $\times7.8$). This explains \texttt{act1}: changing M1--M4 together shifts hits between the
high-momentum group and the others, and the nominal design balances them.

\paragraph{Transverse momentum and charge.}
All $p_t$ bins below $1.5$\,GeV behave alike (worst directions at $\cos=0.89$--$1.00$). The muons
above $1.5$\,GeV ($1.6\%$ of the muons, $23\%$ of the hits) are less sensitive ($0.85$ vs.\
$1.1$--$1.4$ per tolerance unit) and only partly aligned ($\cos=0.66$--$0.72$, Fig.~\ref{fig:pt}).
Positive muons give $64\%$ of the hits; both charges respond to the same changes ($\cos=0.98$), the
$\mu^-$ somewhat more strongly ($1.22$ vs.\ $0.99$).

\begin{figure}[t]
  \centering
  \includegraphics[width=\linewidth]{sim_pz_profiles.png}
  \caption{Simulated hits per $p_z$ range for the nominal design and some designs of
  Table~\ref{tab:concrete}.}
  \label{fig:simpz}
\end{figure}

\begin{figure}[t]
  \centering
  \includegraphics[width=0.9\linewidth]{pz_bins_overview.png}
  \caption{Per $p_z$ bin (surrogate): share of the muons and of the hits, relative sensitivity
  and curvature along the bin's own worst direction.}
  \label{fig:pz_overview}
\end{figure}

\begin{figure}[t]
  \centering
  \includegraphics[width=0.62\linewidth]{pz_bins_alignment.png}
  \caption{Cosine between the worst directions of the $p_z$ bins: the intermediate momenta
  ($30$--$75$\,GeV) want the opposite changes to the high ($>125$\,GeV) and the soft
  ($10$--$30$\,GeV) ones.}
  \label{fig:pz_alignment}
\end{figure}

\begin{figure}[p]
  \centering
  \includegraphics[width=\linewidth]{pz_bins_sensitivity.png}
  \caption{Sensitivity of each $p_z$ bin to each parameter, relative to the bin's own hits.
  Left: first order. Right: exact one-at-a-time change at the worse of $\pm10\%$.}
  \label{fig:pz_sensitivity}
\end{figure}

\begin{figure}[t]
  \centering
  \includegraphics[width=0.49\linewidth]{pt_bins_overview.png}\hfill
  \includegraphics[width=0.42\linewidth]{pt_bins_alignment.png}
  \caption{Per $p_t$ bin: overview (left) and alignment of the bins' worst directions (right).}
  \label{fig:pt}
\end{figure}

\clearpage
%==============================================================================
\section{Conclusions}
%==============================================================================

\begin{enumerate}
  \item \textbf{The design is fragile to coordinated deviations.} With every parameter off by up
        to $\varepsilon$, the worst design found (on the surrogate, then simulated) gives
        $\times1.7$ at $\varepsilon=1\%$, $\times3.8$ at $2\%$, $\times9.0$ at $3\%$, $\times132$ at
        $5\%$ and $\times3440$ at $10\%$; these are lower bounds on the true worst case, tight up
        to $3\%$. Moving every parameter by $1.5\%$ with the sign of its gradient already doubles
        the hits.
  \item \textbf{Bounding the worst case needs sub-percent tolerances on a few parameters and a
        few percent on most.} The allocation found for at most $\times3$ holds M7's field to
        $\pm0.5\%$, M2's entrance width and M7's length to $\pm0.55\%$ and M7's exit width to
        $\pm0.7\%$, 14 parameters below $\pm2\%$, median $5\%$; for at most $\times2$, eleven below
        $\pm1\%$ (simulated $\times2.44$ and $\times1.56$). M7, M2 and M3 need the tightest
        tolerances, the gaps between magnets the loosest.
  \item \textbf{Single parameters underestimate the risk.} The worst one at its $\pm10\%$ limit
        (the field of M7) gives $\times2.1$; the entrance width of M2 and the exit width and length
        of M7 give $\times1.65$--$2.0$. The top three changes of the worst direction together
        give $\times7$, the top ten $\times34$: the factors multiply. M7 matters most, then M2,
        M3 and M6.
  \item \textbf{Random deviations are much milder than the worst case.} Independent errors of
        $\pm2\%$ keep the hits within about $\pm20\%$; at $\pm5\%$ most designs stay within
        $\pm50\%$ but some quadruple the hits; at $\pm10\%$ the median design is twice as bad and
        the tail reaches $\times7$--$29$. A specification can therefore be much looser if
        the errors are independent than if it must hold against the worst combination.
  \item \textbf{Different muons pull in opposite directions.} The hits come from muons of
        $30$--$75$\,GeV ($66\%$) and above $150$\,GeV ($24\%$). Every harmful direction near the
        nominal design acts through the first group, while the high-momentum hits fall; the
        reverse direction raises the soft and high-momentum hits instead. The dominant
        active-subspace direction (the lengths and fields of M1--M4) trades hits between the
        groups and is harmful on both sides.
  \item \textbf{The surrogate is quantitative within the tolerance region.} It reproduces the
        one-at-a-time effects to $4\%$, the walks within one tolerance unit to $1$--$2\sigma$ and
        the worst designs within $\pm1$--$3\%$ to $3$--$5\%$. It overestimates the worst designs
        at $\pm5\%$ ($\times1.8$) and $\pm10\%$ ($\times30$) and is $\sim13\%$ uncertain far from
        the nominal design, so every number used for a decision here is simulated.
\end{enumerate}

\section*{Limitations}
\begin{itemize}
  \item Uniform fields per magnet block, no field maps.
  \item Muons come from the generative flow, so the absolute numbers depend on how well the flow
        reproduces the muon spectrum.
  \item The surrogate's errors are the scatter of its members: they miss a bias the members
        share, such as the one far from the nominal design (Sec.~\ref{sec:validation}). Only the
        top Hessian eigenvector is individually determined.
  \item The random-tolerance distributions rest on eight simulated designs per level and on the
        surrogate for their tails.
  \item The worst designs within a tolerance are found on the surrogate: their simulated values
        bound the true worst case from below, tightly up to $\pm3\%$, possibly loosely at $\pm5\%$
        and $\pm10\%$, where the surrogate is unreliable. The tolerance allocation is a greedy
        heuristic, not an optimum.
\end{itemize}

\clearpage
\appendix
\section{Free parameters}
\begin{table}[h]
\centering
\caption{Free parameters and nominal values (lengths in cm, $B$ in T).}
\label{tab:params}
\small
\begin{tabular}{lrrrrrrrrr}
\toprule
magnet & \texttt{Z\_gap} & \texttt{Z\_len} & \texttt{dXIn} & \texttt{dXOut} & \texttt{gapIn} & \texttt{gapOut} & \texttt{dXYokeIn} & \texttt{dXYokeOut} & $B$ \\
\midrule
M1 &  --   & 151.5 & 63.12 & --    & 8.149 & --    & 64.12 & --     & 1.9 \\
M2 & 15.0  & 224.9 & 68.04 & --    & 8.431 & --    & 70.62 & --     & 1.9 \\
M3 & 16.12 & 225.0 & 64.81 & --    & 8.986 & --    & 85.88 & --     & 1.9 \\
M4 & 15.0  & 224.0 & 39.02 & --    & 8.327 & --    & 97.53 & --     & 1.9 \\
M5 & 15.0  & 120.4 & --    & --    & --    & --    & --    & --     & -- \\
M6 & 15.0  & 168.0 & 20.0  & 30.0  & 24.92 & 9.388 & 97.18 & 102.7  & $-1.2$ \\
M7 & 15.0  & 193.1 & 37.95 & 67.81 & 8.000 & 8.000 & 124.8 & 94.92  & $-1.9$ \\
\bottomrule
\end{tabular}
\end{table}

\section{Integrating the surrogate over the muons}
\label{app:quadrature}

With $A(z)$ the kinematic envelope of the sampled muons and $N=10^9$,
\begin{equation}
  H(u) = \frac{N}{\int\!\mathrm{d}^2z\,\mathcal{N}(z)A(z)}\sum_{q=\pm1}P(q)\int\!\mathrm{d}^2z\,
  \mathcal{N}(z)\,A(z)\int_0^{2\pi}\!\frac{\mathrm{d}\varphi}{2\pi}\,p\big(\text{hit}\mid u, x(z,\varphi,q)\big),
\end{equation}
where $z$ is the flow's 2-d momentum latent, $\varphi$ the azimuth (uniform) and $q$ the charge
(with the data's $\mu^+$ fraction). The derivatives use $801^2$ latent points on $[-7,7]^2$ with
Gaussian weights, 64 azimuths and both charges; everything else $201^2\times64\times2$ points, with
the trunk features cached. The grid is converged to $10^{-5}$ from about $10^7$ points ($3\times10^{-4}$ at $10^6$), while
$10^9$ sampled muons scatter by up to $2\%$ in $H_0$ and $5\%$ above $300$\,GeV, and quasi-Monte
Carlo helps only by a factor two (Fig.~\ref{fig:quadrature}). Importance-sampling the simulator
at a raised temperature does not help either: the hits that matter are at intermediate momenta,
of which a hotter sample holds fewer.

\begin{figure}[h]
  \centering
  \includegraphics[width=\linewidth]{study_quadrature_convergence.png}\\[1ex]
  \includegraphics[width=0.85\linewidth]{study_quadrature.png}
  \caption{Top: Monte Carlo estimates (six seeds) against the grid at increasing resolution, for
  the nominal hits, the hits above $300$\,GeV, the field of M7 at $-10\%$ and one tolerance unit
  along the worst direction.
  Bottom: relative error against the finest grid at equal numbers of points, for Monte Carlo,
  randomised quasi-Monte Carlo (Sobol) and the grid.}
  \label{fig:quadrature}
\end{figure}

\section{How the surrogate was chosen}
\label{app:study}

The choices come from a study of the surrogates (\texttt{surrogate\_study/}) on four toy problems
with exact answers (a binary outcome per sample with known probability; 10--43 parameters) and on
the shield against the reference designs (Fig.~\ref{fig:study}).
\begin{itemize}
  \item \textbf{Modelling each muon is what makes it work.} A response surface fitted to the same
        simulations without the muons needs about ten times more simulation for the same gradient.
  \item \textbf{Training data.} Many designs with few muons each beat few designs with many; all
        designs at once beat a rolling buffer; the designs must cover the region where the
        derivatives are read. A hotter momentum temperature ($T=2$) is much better than $T=1$ or
        $1.5$, and $T=3$ no better than $T=2$. A tailored training proposal (the provably best one
        for the error on $H$) would give at most $1.7\times$ the effective training muons of
        $T=2$, so the temperature is kept. Adding designs far from the nominal one helped only
        with the adaptive pool: the uniform pool alone did not lower the box error on the nominal
        muons, with the adaptive pool it fell by a fifth, to $13\%$ (four seeds each, scored on
        the same references).
  \item \textbf{Architecture.} The DeepONet is the most reliable for the hit rate. At the final
        scale ($2\times10^8$ rows) capacity matters (128 units, 32 features), and averaging seeds
        removes most of the scatter of the gradient. Reduced Taylor surrogates (a per-muon
        quadratic on a low-dimensional subspace), the best Hessians on the toys, are worse here
        (one-at-a-time secants $6$--$10\%$, curvatures $18$--$25\%$). An MLP on $[u,x]$ predicts
        somewhat better (box $11\%$ instead of $16\%$, curvatures $10\%$ instead of $12\%$) and
        agrees with the DeepONet's gradient, top eigenvector and worst direction to
        $\cos=0.99$, so the conclusions do not depend on the choice; the DeepONet is kept because
        its separable form allows cached trunk features and exact derivatives.
\end{itemize}

\begin{figure}[h]
  \centering
  \includegraphics[width=0.53\linewidth]{study_toy_budget.png}\hfill
  \includegraphics[width=0.45\linewidth]{study_shield_budget_final.png}
  \caption{Left: toy problems, errors of the hit rate, gradient and Hessian, and the share of the
  true worst case found, against the simulation budget, per architecture. Right: shield, errors
  on the reference designs against the budget (single models); solid red: the final surrogate,
  dotted: simulation noise.}
  \label{fig:study}
\end{figure}

\end{document}
